\aufzaehlungdrei{Untuk setiap
\mathl{(t,u)}{} yang tetap, terdapat suatu pemetaan linear karena $v$ dan $w$ memasuki suku ketiga secara linear. Jika di sebelah kiri kita mulai dengan
\mathl{(t,u,w)}{}, titik ini dipetakan ke
\mathl{(t,u,0,w)}{} dan kemudian kembali ke
\mathl{(t,u,w)}{,} sehingga memang terdapat suatu pemisahan.
}{Pemetaan singgung dari
\maabbeledisp {u} { I } { I \times \R^n
} { t } { (t,u(t))
} {,}
adalah
\maabbeledisp {} {TI = I \times \R } { T \left( I \times \R^n  \right) = I \times \R^n \times \R \times \R^n
} { (t,a) } { (t,u(t), a, a u'(t) )
} {.}
Oleh karena itu,
\zusatzklammer {turunan vertikal ke arah $\partial$ bersesuaian dengan

\mavergleichskettek
{\vergleichskettek
{ a
}
{ = }{ 1
}
{  }{
}
{    }{
}
{   }{
}
}
{}{}{}} {} {}

\mavergleichskettedisp
{\vergleichskette
{  \nabla_\partial (u)
}
{ =} { \pi_{\text{vert} }  (t,u(t), 1,  u'(t) )
}
{ =} { (t,u(t), -  F \left(  t, u_1 (t) , \ldots , u_n(t)  \right)  +  u'(t)
}
{ } {
}
{ } {
}
}
{}{}{.}
}{Penampang $u$ merupakan
\definitionsverweis {penampang horizontal}{}{}
jika dan hanya jika turunan vertikalnya lenyap, yaitu jika

\mavergleichskettedisp
{\vergleichskette
{  -  F \left(  t, u_1 (t) , \ldots , u_n(t)  \right)  +  u'(t)
}
{ =} { 0
}
{ } {
}
{ } {
}
{ } {
}
}
{}{}{}
atau, secara ekuivalen,

\mavergleichskettedisp
{\vergleichskette
{ u'(t)
}
{ =} {  F \left(  t, u_1 (t) , \ldots , u_n(t)  \right)
}
{ } {
}
{ } {
}
{ } {
}
}
{}{}{}
untuk semua

\mavergleichskette
{\vergleichskette
{ t
}
{ \in }{  I
}
{  }{
}
{    }{
}
{   }{
}
}
{}{}{.}
Ini tepat berarti bahwa kurva tersebut merupakan solusi sistem persamaan diferensial.
}

 [[Kategorie:Latexseite]]
